\documentclass[11pt,a4paper]{article}
\usepackage[spanish,es-noquoting]{babel}
\usepackage[top=2cm, bottom=2cm, left=2.2cm, right=2.2cm]{geometry}
\usepackage{fontspec}
\usepackage{amsmath,amssymb}
\usepackage{enumitem}
\usepackage{tcolorbox}
\usepackage{microtype}

\setmainfont{Calibri}
\linespread{1.15}

\tcbset{
    colback=white,
    colframe=black,
    boxrule=0.8pt,
    arc=2mm,
    left=6pt, right=6pt, top=4pt, bottom=4pt
}

\begin{document}
\pagestyle{empty}

\begin{center}
    {\LARGE \textbf{Resolución Detallada: Examen de Cinemática Vectorial}}\\[4pt]
    {\large \textbf{Física 1.º de Bachillerato}}
\end{center}

\vspace{0.3cm}
\noindent\textbf{Enunciado:} Una partícula efectúa un movimiento cuya ecuación vectorial está determinada por:
\[
\vec{r}(t) = 
\begin{cases}
x(t) = 3t + 6 \\
y(t) = 9t^2 - 15
\end{cases}
\quad (\text{expresada en metros y segundos en el S.I.})
\]

\noindent Es decir: $\vec{r}(t) = (3t + 6)\,\vec{i} + (9t^2 - 15)\,\vec{j}\text{ (m)}$.

\vspace{0.2cm}
\begin{enumerate}[label=\textbf{\alph*)}, leftmargin=2em, itemsep=8pt]
    \item \textbf{El vector de posición en el instante inicial ($t = 0\text{ s}$):}
    Sustituimos $t = 0$ en la ecuación vectorial:
    \[
    \vec{r}(0) = [3(0) + 6]\,\vec{i} + [9(0)^2 - 15]\,\vec{j} = \mathbf{6\,\vec{i} - 15\,\vec{j}\text{ m}} \quad \text{o el punto } (6, -15)\text{ m}
    \]

    \item \textbf{Vector desplazamiento en los dos primeros segundos ($\Delta t = [0, 2]\text{ s}$):}
    \begin{itemize}[leftmargin=1.5em, itemsep=3pt]
        \item Posición en $t = 2\text{ s}$:
        \[
        \vec{r}(2) = [3(2) + 6]\,\vec{i} + [9(2)^2 - 15]\,\vec{j} = (6 + 6)\,\vec{i} + (36 - 15)\,\vec{j} = 12\,\vec{i} + 21\,\vec{j}\text{ m}
        \]
        \item Desplazamiento ($\Delta\vec{r}_{0\to 2} = \vec{r}(2) - \vec{r}(0)$):
        \[
        \Delta\vec{r} = (12 - 6)\,\vec{i} + [21 - (-15)]\,\vec{j} = \mathbf{6\,\vec{i} + 36\,\vec{j}\text{ m}}
        \]
        \item Módulo: $|\Delta\vec{r}| = \sqrt{6^2 + 36^2} = \sqrt{36 + 1296} = \sqrt{1332} = 6\sqrt{37} \approx \mathbf{36{,}50\text{ m}}$.
    \end{itemize}

    \item \textbf{El vector velocidad media en los dos primeros segundos:}
    \[
    \vec{v}_m = \frac{\Delta\vec{r}}{\Delta t} = \frac{6\,\vec{i} + 36\,\vec{j}\text{ m}}{2\text{ s} - 0\text{ s}} = \mathbf{3\,\vec{i} + 18\,\vec{j}\text{ m/s}}
    \]
    \[
    |\vec{v}_m| = \sqrt{3^2 + 18^2} = \sqrt{9 + 324} = \sqrt{333} = 3\sqrt{37} \approx \mathbf{18{,}25\text{ m/s}}
    \]

    \item \textbf{La ecuación de la trayectoria, $y(x)$:}
    \begin{itemize}[leftmargin=1.5em, itemsep=3pt]
        \item Despejamos el parámetro $t$ de la componente $x$:
        \[
        x = 3t + 6 \implies 3t = x - 6 \implies t = \frac{x - 6}{3}
        \]
        \item Sustituimos en $y(t)$:
        \[
        y(x) = 9\left(\frac{x - 6}{3}\right)^2 - 15 = 9 \cdot \frac{(x - 6)^2}{9} - 15 = (x - 6)^2 - 15
        \]
        \item Desarrollando el binomio al cuadrado:
        \[
        y(x) = x^2 - 12x + 36 - 15 \implies \mathbf{y(x) = x^2 - 12x + 21}
        \]
        \textit{(Se trata de una trayectoria parabólica con ramas hacia arriba)}.
    \end{itemize}

    \item \textbf{El vector velocidad y su módulo en el instante $t = 1\text{ s}$:}
    Derivamos el vector de posición respecto al tiempo:
    \[
    \vec{v}(t) = \frac{d\vec{r}(t)}{dt} = \frac{d(3t + 6)}{dt}\,\vec{i} + \frac{d(9t^2 - 15)}{dt}\,\vec{j} = 3\,\vec{i} + 18t\,\vec{j}\text{ m/s}
    \]
    Para $t = 1\text{ s}$:
    \[
    \vec{v}(1) = 3\,\vec{i} + 18(1)\,\vec{j} = \mathbf{3\,\vec{i} + 18\,\vec{j}\text{ m/s}}
    \]
    \[
    v(1) = |\vec{v}(1)| = \sqrt{3^2 + 18^2} = \sqrt{333} = 3\sqrt{37} \approx \mathbf{18{,}25\text{ m/s}}
    \]

    \item \textbf{El vector aceleración y su módulo en ese instante ($t = 1\text{ s}$):}
    Derivamos el vector velocidad:
    \[
    \vec{a}(t) = \frac{d\vec{v}(t)}{dt} = 0\,\vec{i} + 18\,\vec{j} = \mathbf{18\,\vec{j}\text{ m/s}^2}
    \]
    En $t = 1\text{ s}$ (y en cualquier instante, ya que es constante):
    \[
    \vec{a}(1) = \mathbf{18\,\vec{j}\text{ m/s}^2} \qquad a = |\vec{a}| = \mathbf{18\text{ m/s}^2}
    \]

    \item \textbf{La aceleración tangencial ($a_t$) en ese instante ($t = 1\text{ s}$):}
    Proyección del vector aceleración en la dirección de la velocidad (mediante producto escalar):
    \[
    a_t = \frac{\vec{v} \cdot \vec{a}}{|\vec{v}|} = \frac{(3)(0) + (18)(18)}{\sqrt{333}} = \frac{324}{\sqrt{333}} = \frac{324}{3\sqrt{37}} = \frac{108}{\sqrt{37}} \approx \mathbf{17{,}75\text{ m/s}^2}
    \]

    \item \textbf{La aceleración normal ($a_n$) en ese instante ($t = 1\text{ s}$):}
    Como $\vec{a} = \vec{a}_t + \vec{a}_n$ y forman un triángulo rectángulo ($a^2 = a_t^2 + a_n^2$):
    \begin{align*}
    a_n &= \sqrt{a^2 - a_t^2} = \sqrt{18^2 - \left(\frac{108}{\sqrt{37}}\right)^2} = \sqrt{324 - \frac{11664}{37}} \\
    &= \sqrt{\frac{324}{37}} = \frac{18}{\sqrt{37}} \approx \mathbf{2{,}96\text{ m/s}^2}
    \end{align*}
    \textit{Comprobación alternativa por producto vectorial:}
    \[
    a_n = \frac{|\vec{v} \times \vec{a}|}{|\vec{v}|} = \frac{|(3)(18) - (18)(0)|}{\sqrt{333}} = \frac{54}{3\sqrt{37}} = \frac{18}{\sqrt{37}} \approx 2{,}96\text{ m/s}^2
    \]

    \item \textbf{El radio de curvatura ($\rho$) en ese instante ($t = 1\text{ s}$):}
    A partir de la aceleración centrípeta/normal ($a_n = \frac{v^2}{\rho}$):
    \[
    \rho = \frac{v^2}{a_n} = \frac{333}{\frac{18}{\sqrt{37}}} = \frac{333 \cdot \sqrt{37}}{18} = \frac{37\sqrt{37}}{2} \approx \mathbf{112{,}53\text{ m}}
    \]
\end{enumerate}

\end{document}

